% Integrator Concentration in Building Automation Networks
% Supply-chain attack surfaces from shared Foreign Device endpoints

\documentclass[11pt,a4paper]{article}

\usepackage[utf8]{inputenc}
\usepackage[T1]{fontenc}
\usepackage{lmodern}
\usepackage{microtype}
\usepackage{amsmath,amssymb,amsfonts}
\usepackage[numbers,sort&compress]{natbib}
\bibliographystyle{unsrtnat}
\usepackage{hyperref}
\usepackage{url}
\usepackage{booktabs}
\usepackage{graphicx}
\usepackage[margin=1in]{geometry}
\usepackage{xcolor}

\hypersetup{
  pdfauthor={Cameron Warren},
  pdftitle={Integrator Concentration in Building Automation Networks},
  colorlinks=true,
  linkcolor=blue,
  citecolor=blue
}

\title{\textbf{Integrator Concentration}: Supply-Chain Attack Surfaces\\
       in Cloud-Bridged Building Automation Networks}
\author{Cameron Warren\\
  \texttt{cwarre33@uncc.edu}\\[0.5em]
  Department of Computer Science\\
  University of North Carolina at Charlotte
}
\date{\today}

\begin{document}

\maketitle

\begin{abstract}
We discover a previously uncharacterized pattern in building automation 
infrastructure security: \textbf{integrator concentration}, where a single 
third-party endpoint maintains simultaneous Foreign Device registrations 
with multiple geographically separated building control networks. This 
concentration creates a supply-chain attack surface---compromise of one 
monitoring station yields access to multiple critical facilities. Through 
passive analysis of internet-exposed BACnet Broadcast Management Devices, 
we identify and characterize one such concentrator bridging at least two 
distinct building networks over a 47-day observation window. We discuss 
responsible disclosure implications and propose detection methodologies 
for identifying concentrator infrastructure at scale.
\end{abstract}

\section{Introduction}
\label{sec:intro}

Supply-chain attacks have emerged as a dominant pattern in cyber-physical 
system compromises: rather than targeting hardened endpoints directly, 
adversaries compromise shared infrastructure (software updates, managed 
services, third-party integrations) to gain access to multiple downstream 
targets.

Building automation systems have been considered relatively isolated from 
threat vectors that plague networked industrial control systems. We 
discover this isolation assumption may be incorrect.

\subsection{The Concentration Phenomenon}

Through analysis of internet-exposed BACnet BBMDs (Building Broadcast 
Management Devices), we identify a pattern where multiple building networks 
maintain independent tunnels to a shared external endpoint---suggesting 
third-party integrator management infrastructure bridging client facilities.

This pattern enables:
\begin{itemize}
  \item Horizontal movement between facility networks via integrator infrastructure
  \item Scale advantages for adversaries (one compromise, multiple targets)
  \item Attribution challenges (attack appears to originate from trusted third-party)
\end{itemize}

\subsection{Key Finding}

Two distinct BBMDs, separated by approximately 4,000 km in different 
municipalities, maintain coordinated Foreign Device registrations with 
the same public IP address (216.67.73.166) over a 47-day observation 
window.

\subsubsection{Contributions}

\begin{enumerate}
  \item \textbf{Formalization of Integrator Concentration}: We define and 
    characterize a novel supply-chain attack surface in building automation.
    
  \item \textbf{Empirical Discovery}: We identify and document a real 
    concentrator endpoint serving at least two distinct building networks 
    over an extended period.
    
  \item \textbf{Detection Framework}: We propose passive OSINT methods 
    for identifying concentrator infrastructure at scale.
    
  \item \textbf{Disclosure Study}: We discuss responsible disclosure 
    challenges when both integrator and building operator are affected 
    parties.
\end{enumerate}

\section{Background}
\label{sec:background}

\subsection{BACnet Foreign Device Registration}

BACnet supports cross-subnet communication through Foreign Device 
registration: a device on one subnet registers with a BBMD on another, 
creating a tunnel for forwarded broadcasts. Architecturally, this enables:

\begin{itemize}
  \item Multi-building campus networks (intended use)
  \item Remote monitoring by third-party integrators (accepted use)
  \item Cloud-managed building automation (emerging pattern)
\end{itemize}

Standard security assumptions: Foreign Devices are within the same 
organization's control.

\subsection{Supply-Chain Attacks in Industrial Control}

Prior research documents:
\begin{itemize}
  \item Software supply chains (compromised updates affecting multiple facilities)
  \item MSP (Managed Service Provider) compromises cascading to clients
  \item Vendor credential breaches exposing customer networks
\end{itemize}

\textbf{Gap}: No prior work has characterized third-party infrastructure 
concentration in building automation specifically.

\section{Discovery and Methodology}
\label{sec:methodology}

\subsection{Data Source}

Longitudinal Shodan dataset: BBMDs with BACnet/WP fingerprints sampled 
2026-03-04 to 2026-04-20, with extracted Foreign Device Tables.

\subsection{Concentration Detection}

We define concentrator indicators:

\begin{enumerate}
  \item \textbf{Multi-host FD registration}: Same FD IP across multiple BBMDs
  \item \textbf{Temporal overlap}: Simultaneous or overlapping registrations
  \item \textbf{Geographic separation}: BBMDs in distinct locations
  \item \textbf{Behavioral similarity}: Similar port/TTL patterns
\end{enumerate}

\subsection{Confirmed Concentrator: 216.67.73.166}

\begin{table}[h]
\centering
\begin{tabular}{ll}
\toprule
\textbf{Attribute} & \textbf{Observation} \\
\midrule
Connected BBMDs & Minimum 2 confirmed \\
BBMD 1 & 66.58.248.125 (United States) \\
BBMD 2 & 24.237.132.230 (United States, different region) \\
Registration duration & 47 days continuous \\
Port behavior & 13 rotating source ports each \\
Last seen & 2026-04-20 (both BBMDs) \\
\bottomrule
\end{tabular}
\caption{Confirmed Concentrator Characteristics}
\label{tab:concentrator}
\end{table}

\subsection{Passive Attribution}

Without active connection to the concentrator endpoint, we infer likely 
integrator nature from:

\begin{itemize}
  \item Rotating-port behavior (suggesting NAT traversal, common in 
    consumer-grade integrator setups)
  \item Persistence over 47 days (not transient, intentional service)
  \item Geographic distribution (not colocated, infrastructure service)
\end{itemize}

\section{Attack Surface Analysis}
\label{sec:attacks}

\subsection{Threat Model}

\textbf{Adversary capability}: Compromise of integrator monitoring station.

\textbf{Attack path}:
\begin{enumerate}
  \item Gain access to integrator endpoint through phishing, credential 
    theft, or infrastructure vulnerability
  \item Leverage existing Foreign Device registrations as persistent 
    tunnels into client building networks
  \item Issue BACnet commands from trusted FD position
\end{enumerate}

\textbf{Advantages over direct facility targeting}:
\begin{itemize}
  \item Single compromise, multiple facilities
  \item Appears as legitimate integrator traffic
  \item Bypasses facility-level perimeter defenses
\end{itemize}

\subsection{Impact Scenarios}

Building automation systems control:
\begin{itemize}
  \item HVAC (climate, air handling, temperature)
  \item Access control (doors, elevators)
  \item Fire safety systems
  \item Critical infrastructure cooling
\end{itemize}

Compromise enables:
\begin{itemize}
  \item Climate manipulation affecting equipment/occupants
  \item Access bypass at critical facilities
  \item Ransomware pivot to IT networks via HVAC bridge
\end{itemize}

\section{Detection at Scale}
\label{sec:detection}

\subsection{Large-Scale Concentrator Discovery}

Proposed methodology for passive detection:

\begin{enumerate}
  \item Harvest full Shodan BACnet dataset ($\sim$40,000 BBMDs indexed)
  \item Extract all Foreign Device Table entries
  \item Group by FD IP, counting distinct BBMD registrations
  \item Filter for multiple BBMDs per FD with geographic separation
  \item Validate persistence over time window
\end{enumerate}

\textbf{Expected yield}: Based on our 17-BBMD sample showing 1 concentrator, 
full dataset may reveal 50--100 additional concentrator endpoints.

\subsection{False Positive Mitigation}

\begin{itemize}
  \item \textbf{Exclude colocated BBMDs}: Same datacenter = different threat model
  \item \textbf{Require temporal overlap}: Not sequential use
  \item \textbf{Behavioral similarity}: Consistent port/TTL patterns
\end{itemize}

\section{Responsible Disclosure Challenges}
\label{sec:disclosure}

\subsection{Multi-Party Affected}

Concentrator findings implicate multiple parties:

\begin{itemize}
  \item Building operators (facilities with exposed BBMDs)
  \item Integrator (owner of concentrator endpoint)
  \item Potentially building occupants/residents
\end{itemize}

Standard disclosure assumes single responsible party.

\subsection{Attribution Barriers}

Identifying integrators from passive data requires:
\begin{itemize}
  \item IP-to-ASN mapping
  \item WHOIS lookup (often privacy-protected)
  \item Business registration search
\end{itemize}

Failure modes: privacy-protected WHOIS, offshore registration, 
home-office integrators.

\subsection{Proposed Approach}

\begin{enumerate}
  \item Attempt integrator attribution first
  \item If successful: disclose to integrator as primary responsible party
  \item Notify building operators with guidance on FD registration audit
  \item Coordinate remediation timing
\end{enumerate}

\section{Conclusion}

Integrator Concentration represents a novel supply-chain attack surface 
in building automation infrastructure. Our empirical discovery of a 
confirmed concentrator bridging multiple facilities demonstrates the 
pattern is not theoretical---it is actively deployed and discoverable 
through passive OSINT methods.

The concentration phenomenon suggests building automation security 
assessments must extend beyond facility boundaries to include third-party 
management infrastructure that provides bridged access.

\section*{Acknowledgments}

This research was conducted using only publicly indexed Shodan data.
All findings have been reported to US-CERT for coordination with 
affected parties.

\end{document}
